\documentclass[11pt]{article}

\usepackage[T1]{fontenc}
\usepackage[utf8]{inputenc}
\usepackage{lmodern}
\usepackage{microtype}
\usepackage[a4paper,margin=27mm]{geometry}
\usepackage{amsmath,amssymb,amsthm,mathtools}
\usepackage{enumitem}
\usepackage{booktabs}
\usepackage{hyperref}
\usepackage[nameinlink,capitalise,noabbrev]{cleveref}

\hypersetup{
  colorlinks=true,
  linkcolor=blue,
  citecolor=blue,
  urlcolor=blue,
  pdftitle={Modal Triplet Theory and History-Dependent Stochastic Processes},
  pdfauthor={Peter Nero}
}

\newtheorem{theorem}{Theorem}[section]
\newtheorem{proposition}[theorem]{Proposition}
\newtheorem{lemma}[theorem]{Lemma}
\newtheorem{corollary}[theorem]{Corollary}
\theoremstyle{definition}
\newtheorem{definition}[theorem]{Definition}
\newtheorem{example}[theorem]{Example}
\theoremstyle{remark}
\newtheorem{remark}[theorem]{Remark}

\newcommand{\R}{\mathbb{R}}
\newcommand{\N}{\mathbb{N}}
\newcommand{\PP}{\mathbb{P}}
\newcommand{\EE}{\mathbb{E}}
\newcommand{\cB}{\mathcal{B}}
\newcommand{\cF}{\mathcal{F}}
\newcommand{\cH}{\mathcal{H}}
\newcommand{\cA}{\mathcal{A}}
\newcommand{\cZ}{\mathcal{Z}}
\newcommand{\TV}{\mathrm{TV}}

\title{Modal Triplet Theory and History-Dependent Stochastic Processes:\\
Fixed-State Indivisibility, Markov Order, and the Quantum Boundary}
\author{Peter Nero}
\date{September 2026 (v3)}

\begin{document}
\maketitle

\begin{abstract}
This paper gives a measure-theoretic account of stochastic processes
obtained by projecting an upper Modal Triplet Theory (MTT) evolution.  Given a
selected probability law on upper initial data, measurable upper dynamics, and
an observable map, the projected histories carry a unique path-space law.
Regular conditional kernels then exist on standard Borel state spaces.  The
observed process may fail to be Markovian even when the upper process is
Markovian or deterministic.  We distinguish four statements that were
previously conflated: failure of the Markov property on the observed state
space, absence of finite Markov order, failure of a chosen
Chapman--Kolmogorov reduction, and quantum CP-indivisibility.  Every classical
path law still factorizes sequentially through its history-dependent
conditional kernels and becomes Markovian on a history state space.  An exact
binary hidden-state example shows that projection can create infinite observed
Markov order while the memory remains compressible to one posterior scalar.
Exponential action weights define valid kernels when their action and reference
measure are supplied, but writing a positive kernel in Gibbs form is a
parameterization, not a first-principles prediction.  Finally, a commutative
path-space algebra does not derive noncommutative quantum observables,
process tensors, Born probabilities, or quantum instruments.  This limitation
does not reopen the separately established canonical operational q79 model.
We explain its relation to coherent history witnesses, frozen finite
hierarchical-equations-of-motion calculations, and bounded optical readout.
A charged molecular line and conditional geometric repair provide further
comparison structures, but neither an abstract intertwiner nor a fitted
response selects a physical q79--FMO correspondence.  The remaining physical
obligations are a same-source preparation and intervention model, controlled
observable transport, calibration, and an independently fixed clock.
\end{abstract}

\section*{Revision note: v3}
\paragraph{Supersedes.}
Version 3 supersedes v2 (July 2026; Zenodo record 21665996). The earlier
revision note below is retained as a record of that correction, not as the
current status of every quantum-source question.
\paragraph{Reason.}
The classical account predates the canonical operational q79 source and the
subsequent BEQ coherent-history, readout, charged-line, and repair results.
Its blanket quantum-source boundary and older corpus ledger were stale.
\paragraph{Resolution.}
The new contextual sections distinguish exact constructions, finite numerical
profiles, excluded realizations, conditional comparisons, and physical source
selection. They explain all 48 assigned frozen results, including later
refinements inside the evolving candidate audit. The transported-metric
semigroup estimate now states its common-norm assumptions explicitly.
\paragraph{Retained.}
The classical projection and Markovization proofs, the posterior-memory
example, and measurement as ordinary physical intervention are retained.
No finite q79 or canonical operational quantum conclusion is reopened.
\paragraph{Open boundary.}
The physical molecular correspondence, selected nonflat endpoint, detector
and source uncertainty, bounded bath-active realization, and dimensionful
clock remain separate obligations. Pure positive repair is not the selected
physical signed action. Reading frozen certificates is not an independent
scientific rerun.

\section*{Revision note}

\paragraph{Supersedes.}
Version 2 supersedes the January 2026 release associated with Zenodo record
18254863.

\paragraph{Reason.}
The former version called the construction complete and first-principles while
assuming its probability source.  It also inferred indivisibility from history
dependence, treated exact realization of a prescribed kernel as an MTT
prediction, and used a commutative GNS representation as if it supplied
general quantum instruments.

\paragraph{Resolution.}
This version replaces those claims by a conditional projection theorem,
separates Markov order from classical and quantum divisibility, proves
history-space Markovization, and distinguishes a physical intervention from
Bayesian conditioning.  It also replaces the inconsistent product of three
three-manifolds by the current upper-space notation \(Y^4\times X^6\), with
the \(1<2<3\) lanes understood as operator data on a shared carrier.

\paragraph{Retained content.}
The valid path-space construction, the possibility of projection-induced
memory, normalized history-dependent kernels, and ordinary conditional
probability are retained with explicit hypotheses.

\paragraph{Remaining boundary.}
MTT must still select an upper capture or preparation law, derive any proposed
action-weight kernel from the same branch, rule out finite sufficient-state
compression when infinite memory is claimed, and provide a noncommutative
observable and instrument source before quantum-process conclusions follow.

\tableofcontents

\section{Introduction: what projection can and cannot do}

The central idea of this paper is simple.  A lower observer may see only a
coarse variable \(q_k\), while an upper state \(z_k\) contains additional
information.  Even if the upper evolution is deterministic, two upper states
with the same present value of \(q_k\) can have different futures.  An ensemble
of such upper states therefore produces an observed process whose prediction
depends on more than the present observed value.

This is a standard and important mechanism.  Hidden Markov models, coarse
graining, and open-system reductions all exhibit it.  MTT gives the mechanism
a geometric interpretation: the observed record is a projection of a richer
fixed-point and bundle state.  The geometry can therefore be a source of
memory in the projected description.

The mechanism alone, however, does not supply probabilities.  A deterministic
map and a projection produce a single history from a single initial state.
To obtain a probability law one must additionally specify a distribution over
initial states, a stochastic upper transition rule, or some other selected
capture measure.  This distinction is the main logical boundary of the
revised paper.

\subsection{Four notions that must be kept separate}

The word ``indivisible'' is used in several inequivalent ways.  We distinguish:

\begin{enumerate}[label=(\roman*)]
\item failure of the one-step Markov property on the observed state space
  \(Q\);
\item absence of every finite Markov order on \(Q\);
\item failure of a specified two-time family to satisfy a
  Chapman--Kolmogorov composition law on \(Q\);
\item failure of CP-divisibility for a family of quantum channels.
\end{enumerate}

The first three are classical properties.  The fourth requires a
noncommutative operator system and completely positive maps.  It cannot be
deduced from a classical path law alone.  Modern process-tensor treatments
make this operational distinction explicit
\cite{pollock2018framework,pollock2018markov,rivas2010}.

\subsection{Claim map}

\begin{center}
\begin{tabular}{p{0.43\textwidth}p{0.45\textwidth}}
\toprule
Statement & Status in this paper \\
\midrule
Selected upper law plus dynamics and projection define a path law
& Proved. \\
Projection can destroy the Markov property on \(Q\)
& Proved conditionally and illustrated exactly. \\
History dependence proves absolute indivisibility
& False; replaced by history-space Markovization. \\
Complete-past dependence proves irreducible infinite memory
& False without excluding finite sufficient statistics. \\
\(\exp(-\Delta A/\eta)\) defines a kernel
& Proved when \(\Delta A,\eta\), and a reference measure are supplied. \\
Every positive kernel has an action representation
& True but tautological; it is not a prediction. \\
Classical conditioning is a quantum instrument
& False without a noncommutative source theorem. \\
Canonical operational q79 source
& Established at its finite-symbol, binary, one-anchor tier. \\
Physical q79--FMO preparation and intervention source
& Open; compatibility and finite readout do not select it. \\
\bottomrule
\end{tabular}
\end{center}

\section{The corrected MTT input}

\subsection{Upper geometry}

The stochastic theorem does not require a literal product of three independent
three-manifolds.  We use the current branch notation
\[
M=Y^4\times X^6,
\]
where \(Y^4\) is the external spacetime and \(X^6\) is the selected compact
internal geometry when such a branch has been supplied.  The Circle--Lens--Nil
\(1<2<3\) structure is treated as a rank or operator filtration on a shared
carrier.  In particular, its lane count is not added as
\(4+3+3+3\), which would be thirteen rather than ten.

For the probability theory below, all geometric and field data are collected
into an upper state space \((\cZ,\cB(\cZ))\).  This may include fields on
\(Y^4\times X^6\), bundle connections, fixed-point data, recorder variables,
and any finite source coordinates.  We assume \(\cZ\) is a standard Borel
space.  Let
\[
T:\cZ\longrightarrow\cZ
\]
be one step of the selected upper evolution, and let
\[
\pi:\cZ\longrightarrow Q
\]
be the measurable observable or recorder map into a standard Borel space
\(Q\).

\subsection{Where probability enters}

A probability measure \(\rho\) on \(\cZ\) is an additional datum.  It may
represent a preparation ensemble, a capture measure on basins, uncertainty in
unobserved upper variables, or a genuinely stochastic source.  Those physical
interpretations are different, but the mathematical construction begins only
after \(\rho\) has been specified.

\begin{theorem}[Projected path law]\label{thm:projected-law}
Let \(\cZ\) and \(Q\) be standard Borel spaces, let \(T:\cZ\to\cZ\) and
\(\pi:\cZ\to Q\) be measurable, and let \(\rho\) be a probability measure on
\(\cZ\).  Define
\[
\Gamma(z)=\bigl(\pi(z),\pi(Tz),\pi(T^2z),\ldots\bigr)\in Q^{\N}.
\]
Then \(\Gamma\) is measurable and
\[
\mu:=\Gamma_{\ast}\rho
\]
is a unique probability measure on
\((Q^{\N},\cB(Q)^{\otimes\N})\).  Its coordinate process
\(Q_k(\omega)=\omega_k\) is the stochastic shadow of the upper ensemble.
\end{theorem}

\begin{proof}
Every coordinate \(\pi\circ T^k\) is measurable.  The product sigma-algebra is
generated by coordinate cylinders, hence \(\Gamma\) is measurable.  The
pushforward of a probability measure under a measurable map is a probability
measure and is uniquely defined by
\(\mu(C)=\rho(\Gamma^{-1}C)\).
\end{proof}

\begin{remark}
If \(\rho=\delta_{z_0}\), then \(\mu\) is concentrated on one path.  Neither
projection nor fixed-point dynamics turns that Dirac law into a nontrivial
probability distribution.  This is why the capture or preparation law is a
real source obligation rather than a technical detail.
\end{remark}

\subsection{Finite phase recurrence is not irreversible mixing}
\label{sec:recurrence-boundary}
For a closed finite-mode self-adjoint phase generator, matrix coefficients
of unitary evolution are finite sums of oscillatory exponentials. They are
almost periodic, not decaying mixing correlations. With compact resolvent,
the same recurrence conclusion holds for individual vectors by approximation
with finite spectral sums. Compact resolvent alone does not turn these
reversible phases into a physical stochastic driver.

This is not a no-go for finite dissipative models: a supplied Markov or
Lindblad generator can mix, and a positive repair heat semigroup can decay.
Those are different dynamics. Deriving irreversible stochastic behavior
from a closed microscopic phase model requires a controlled continuum,
bath, weak-coupling, or thermodynamic limit with its state, observables,
scaling, and error bounds specified. The finite HEOM comparisons below
are calculations within a supplied bath model, not that derivation
\cite{beq04,beq06}. The selected physical driver construction remains a
B.ACTION.01 obligation, without reopening the canonical operational model.

\section{Conditional kernels and action weights}

\subsection{Construction from kernels}

One may equivalently start with an initial law \(\lambda_0\) on \(Q\) and a
sequence of measurable probability kernels
\[
K_k(\,\cdot\,|h_k),\qquad h_k=(q_0,\ldots,q_k)\in Q^{k+1}.
\]
The next theorem is standard probability theory
\cite{kallenberg2002}.

\begin{theorem}[Ionescu--Tulcea extension]\label{thm:IT}
For every initial probability measure \(\lambda_0\) and sequence of probability
kernels \(K_k\) from \(Q^{k+1}\) to \(Q\), there is a unique probability measure
\(\mu\) on \(Q^{\N}\) whose finite-dimensional distributions satisfy
\[
\begin{split}
&\mu(Q_0\in A_0,\ldots,Q_n\in A_n)\\
&\quad =
\int_{A_0}\lambda_0(dq_0)
\int_{A_1}K_0(dq_1|q_0)\cdots
\int_{A_n}K_{n-1}(dq_n|q_0,\ldots,q_{n-1}).
\end{split}
\]
\end{theorem}

Conversely, because \(Q\) is standard Borel, the law in
\Cref{thm:projected-law} admits regular conditional distributions
\[
K_k(A|Q_0,\ldots,Q_k)
=\mu(Q_{k+1}\in A\mid Q_0,\ldots,Q_k)
\]
for every \(k\), unique up to the usual null sets.  Thus every path law
factorizes sequentially through history-dependent kernels.  A claim that a
path law has no such factorization is therefore incorrect.

\subsection{A selected action kernel}

Let \(\nu\) be a sigma-finite reference measure on \(Q\), let
\(\eta>0\), and let
\[
\Delta A_k:Q^{k+1}\times Q\longrightarrow\R\cup\{+\infty\}
\]
be jointly Borel.  Define
\[
Z_k(h_k)=\int_Q
\exp\!\left[-\frac{\Delta A_k(h_k,q')}{\eta}\right]\nu(dq').
\]

\begin{lemma}[Normalized action kernel]\label{lem:action-kernel}
On the Borel set \(G_k=\{h_k:0<Z_k(h_k)<\infty\}\), the expression
\[
K_k(A|h_k)=\frac{1}{Z_k(h_k)}
\int_A\exp\!\left[-\frac{\Delta A_k(h_k,q')}{\eta}\right]\nu(dq')
\]
is a measurable probability kernel. Extend it outside \(G_k\) by any fixed
probability measure. If every reachable history lies in \(G_k\), this extension
does not change the intended process. Together with an initial law the
extended kernel defines a unique path measure by \Cref{thm:IT}.
\end{lemma}

\begin{proof}
Nonnegativity and normalization are immediate.  Measurability of parameterized
nonnegative integrals follows from joint Borel measurability; division is
measurable because \(Z_k\) is finite and strictly positive.
\end{proof}

This result is exact, but it begins with \(\Delta A_k\), \(\eta\), and \(\nu\).
Calling \(\eta=\hbar\) is a physical identification, not a consequence of
normalization.  Likewise, a Euclidean exponential weight is not the complex
phase \(e^{iS/\hbar}\) and does not by itself produce quantum interference.

\begin{lemma}[Gibbs form is a parameterization]\label{lem:gibbs}
Suppose \(K_k(\cdot|h_k)\) has a strictly positive normalized density
\(k_k(\cdot|h_k)\) with respect to \(\nu\).  For any \(\eta>0\), setting
\[
\Delta A_k(h_k,q')=-\eta\log k_k(q'|h_k)
\]
reproduces \(K_k\) through \Cref{lem:action-kernel}, with \(Z_k=1\).  Adding an
arbitrary history-dependent constant to \(\Delta A_k\) leaves the normalized
kernel unchanged.
\end{lemma}

\begin{proof}
Exponentiation gives \(e^{-\Delta A_k/\eta}=k_k\), and the density is already
normalized.  A history-only constant multiplies numerator and denominator by
the same factor.
\end{proof}

\Cref{lem:gibbs} is useful for fitting or encoding a known kernel.  It cannot
show that MTT geometry selected that kernel.  Exact reconstruction of arbitrary
target kernels demonstrates expressive capacity; prediction requires a rule
that fixes the action before the target data are supplied.

\section{Markov order and fixed-state indivisibility}

\subsection{Precise definitions}

Let \(\cF_k=\sigma(Q_0,\ldots,Q_k)\).

\begin{definition}[Markov order]\label{def:markov-order}
The process has Markov order at most \(r\geq 1\) if, for every \(k\geq r-1\),
there is a kernel \(L_k\) such that
\[
\PP(Q_{k+1}\in A\mid\cF_k)
=L_k(A|Q_{k-r+1},\ldots,Q_k)
\quad\text{a.s.}
\]
for every Borel \(A\).  It has infinite Markov order on \(Q\) if no finite
\(r\) has this property.
\end{definition}

\begin{definition}[Fixed-\(Q\) indivisibility]\label{def:q-indivisible}
For this paper only, a process is \emph{fixed-\(Q\) indivisible} if its complete
finite-dimensional distributions cannot be generated by one-step kernels
depending only on the current value \(q_k\).  This is simply failure of Markov
order one on the chosen observed state space.  It is not quantum
CP-indivisibility.
\end{definition}

\begin{proposition}[Factorization criterion]\label{prop:factorization}
A path law is Markov of order one on \(Q\) if and only if there are kernels
\(M_k(dq_{k+1}|q_k)\) for which every finite-dimensional law factorizes as
\[
\lambda_0(dq_0)\prod_{k=0}^{n-1}M_k(dq_{k+1}|q_k).
\]
Consequently, fixed-\(Q\) indivisibility is exactly a statement about the
chosen state description \(Q\), not about the impossibility of sequential
factorization.
\end{proposition}

\begin{proof}
The forward implication is \Cref{thm:IT} applied to Markov kernels.  For the
reverse implication, disintegration of the displayed law gives a version of
the conditional law of \(Q_{k+1}\) that depends only on \(Q_k\).
\end{proof}

\subsection{Every process is Markov on its history state}

\begin{theorem}[History-state Markovization]\label{thm:history-markov}
Let \(\mu\) be any path law on a standard Borel space \(Q\), with conditional
kernels \(K_k\).  Define
\[
\cH=\bigsqcup_{k\geq0}\{k\}\times Q^{k+1},
\qquad
H_k=(k,Q_0,\ldots,Q_k).
\]
Then \(H_k\) is a first-order Markov process on \(\cH\).  Its transition from
\((k,h_k)\) appends \(q'\) with law \(K_k(dq'|h_k)\).
\end{theorem}

\begin{proof}
Given \(H_k\), the complete observed past is known.  The conditional
distribution of \(H_{k+1}\) is therefore the pushforward of
\(K_k(\cdot|h_k)\) under
\(q'\mapsto(k+1,h_k,q')\), and depends on no earlier history state.
\end{proof}

This theorem is not a trick that removes physically relevant memory.  It tells
us what must be measured: the size and structure of the state required to make
the process Markovian.  A useful non-Markovianity claim should therefore rule
out a small sufficient state, not merely observe that \(q_k\) alone is
insufficient.

\begin{proposition}[Sufficient-state compression]\label{prop:sufficient}
Suppose there is a standard Borel space \(S\), measurable summaries
\(s_k:Q^{k+1}\to S\), update maps
\[
s_{k+1}(h_k,q')=u_k(s_k(h_k),q'),
\]
and kernels \(L_k\) such that
\[
K_k(\cdot|h_k)=L_k(\cdot|q_k,s_k(h_k)).
\]
Then \((Q_k,S_k)\), with \(S_k=s_k(Q_0,\ldots,Q_k)\), is first-order Markov.
Complete-past notation therefore does not prove irreducible or
infinite-dimensional memory.
\end{proposition}

\begin{proof}
The next \(Q\)-value depends only on \((Q_k,S_k)\), and the next summary is a
measurable function of \((S_k,Q_{k+1})\).  Hence the joint next-state law
depends only on the current augmented state.
\end{proof}

The proposition permits any standard Borel \(S\). Calling the compression
finite requires a specified finite-state or finite-dimensional class of
summaries; measurable encoding alone is not a physical memory-size bound.

\section{An exact projection-induced memory example}

\begin{example}[Binary latent source]\label{ex:binary}
Let \(S\in\{0,1\}\) be chosen with equal prior probabilities and remain fixed.
Conditional on \(S\), let \(Q_0,Q_1,\ldots\) be independent with
\[
\PP(Q_k=S|S)=1-\varepsilon,\qquad
\PP(Q_k\neq S|S)=\varepsilon,
\qquad 0<\varepsilon<\frac12.
\]
The upper state includes \(S\); the observer records only \(Q_k\).
\end{example}

For a history \(h_k=(q_0,\ldots,q_k)\), write \(n_1\) and \(n_0\) for the
numbers of ones and zeros.  Bayes' rule gives
\[
p_k:=\PP(S=1|h_k)
=\frac{(1-\varepsilon)^{n_1}\varepsilon^{n_0}}
{(1-\varepsilon)^{n_1}\varepsilon^{n_0}
 +\varepsilon^{n_1}(1-\varepsilon)^{n_0}},
\]
and therefore
\[
\PP(Q_{k+1}=1|h_k)
=\varepsilon+(1-2\varepsilon)p_k.
\]

\begin{proposition}[Infinite observed order, one-scalar memory]\label{prop:binary}
The process in \Cref{ex:binary} has no finite Markov order on \(Q=\{0,1\}\).
Nevertheless, the scalar posterior \(p_k\) is a sufficient state, so the
observed memory is recursively compressible to one real number.
\end{proposition}

\begin{proof}
Fix any proposed order \(r\).  Choose two positive-probability histories with
the same final \(r\) symbols but with sufficiently many additional ones in one
prefix and zeros in the other.  Their values of \(n_1-n_0\), hence their
posteriors \(p_k\), differ.  Their next-symbol probabilities therefore differ,
so no order-\(r\) kernel can represent both histories.

For compression, after observing \(q\in\{0,1\}\), Bayes' rule updates \(p\) by
\[
u(p,q)=
\frac{p(1-\varepsilon)^q\varepsilon^{1-q}}
{p(1-\varepsilon)^q\varepsilon^{1-q}
 +(1-p)\varepsilon^q(1-\varepsilon)^{1-q}}.
\]
The next prediction depends only on \(p_k\), and
\(p_{k+1}=u(p_k,Q_{k+1})\).  Apply \Cref{prop:sufficient}.
\end{proof}

This example captures both the promise and the warning for MTT.  Projection
can genuinely produce observed infinite-order memory.  Yet an upper hidden
coordinate, or even a one-dimensional posterior over it, can restore a Markov
description.  A claim of irreducible MTT memory must therefore prove that no
allowed finite recorder state or sufficient statistic closes the evolution.

\section{The classical--quantum boundary}

\subsection{What classical GNS does}

Let \((\Omega,\cF,\mu)\) be the path space.  Bounded classical observables form
the commutative algebra \(L^\infty(\Omega,\mu)\).  Acting by multiplication on
\(L^2(\Omega,\mu)\), an event \(E\) is represented by the projection
\[
(M_Ef)(\omega)=\mathbf{1}_E(\omega)f(\omega),
\]
and
\[
\mu(E)=\langle \mathbf{1},M_E\mathbf{1}\rangle.
\]
This has the appearance of a Born formula, but all \(M_E\) commute.

\begin{proposition}[Commutative reconstruction limit]\label{prop:commutative}
The GNS representation of a commutative path-space algebra has a commuting
operator image.  It does not, without extra structure, produce incompatible
observables, complex interference amplitudes, a noncommutative state algebra,
or general completely positive instruments.
\end{proposition}

\begin{proof}
If \(ab=ba\) in the algebra, then every representation satisfies
\(\pi(a)\pi(b)=\pi(ab)=\pi(ba)=\pi(b)\pi(a)\).  The missing noncommutative
objects cannot be obtained by relabeling commuting multiplication operators.
\end{proof}

Ordinary sigma-additivity also implies that the probability of a fixed
measurable event is unchanged when a partition is refined around it.  That is
a consistency property of one classical sample space, not by itself the
operational measurement noncontextuality used in quantum-foundational
theorems \cite{spekkens2005}.

\subsection{Why quantum divisibility is separate}

Quantum CP-divisibility concerns a family of channels
\(\Lambda_{t:s}\) for which
\[
\Lambda_{t:r}=\Lambda_{t:s}\Lambda_{s:r}
\]
with the intermediate map completely positive and trace preserving.  A
process tensor goes further by mapping sequences of interventions to outcome
statistics.  Both notions require noncommutative input and output systems and
an intervention calculus \cite{pollock2018framework,pollock2018markov}.

A classical kernel can agree numerically with the probabilities of one chosen
quantum experiment.  Such agreement does not identify its updates with all
quantum instruments or establish CP-divisibility. The classical construction
does not supply those structures. A separate canonical q79 construction now
does so at its declared operational tier, as explained next; extending that
model to a selected molecular preparation and detector is a different task.

\subsection{The established operational q79 source}\label{sec:q79-source}

The frozen source-base result imports a noncommutative model on
\(\cH_\Sigma=L^2(\Sigma,d\mu_h;F_{q79})\), with positive trace-class states
and a bounded completion of the decomposable observable algebra. Its binary
finite-symbol generator is
\[
\overline{\mathcal L}(\rho)=\log(448)
 \bigl(P\rho P+Q\rho Q-\rho\bigr),\qquad Q=I-P.
\]
The construction includes the preparation-ensemble second-moment and ready
Fock finite-horizon structure at the canonical one-anchor tier. These are
not inferred from commuting path events or fitted observed frequencies
\cite{beq46}. The theorem here remains a classical theorem; the existence of
that separate source prevents its limitation from being misreported as a
failure of operational MTT quantum mechanics.

An elementary compatibility model further illustrates the distinction. The
nilpotent matrix \(d=E_{12}\) on a three-dimensional carrier has
\(d^\dagger d+dd^\dagger=\operatorname{diag}(1,1,0)\). Exact changes of basis
give the q79 Hodge split and a route-cycle split. This is an explicit
compatibility square, not a uniqueness theorem selecting a biological
Hamiltonian. In particular the full natural q79 unital dephasing generator
is not the full natural FMO HEOM generator. The source's corresponding
no-go is retained; a future restricted observable comparison has a different
type \cite{beq47}.

\section{Measurement: intervention first, conditioning second}

Measurement has no special metaphysical status here.  It is an ordinary
physical interaction involving a system, an apparatus, and an outcome record.
Probability theory describes two different operations:

\begin{enumerate}[label=(\roman*)]
\item an \emph{intervention} changes the physical transition law because the
  apparatus is coupled to the system;
\item \emph{conditioning} updates the probability law after an outcome has
  been learned.
\end{enumerate}

To represent the first operation classically, let \(a_k\) denote an apparatus
setting and use an intervention-dependent kernel
\[
K_k^{a_k}(dq_{k+1},dy_k|h_k),
\]
where \(y_k\) is the record.  The experimental protocol and these kernels
define a path law.  After observing \(Y_k=y\), the posterior is the ordinary
conditional measure
\[
\mu^{a,y}(B)=\frac{\mu^a(B\cap\{Y_k=y\})}
{\mu^a(Y_k=y)}
\]
when the denominator is nonzero.

Replacing the physical intervention by posterior conditioning loses
disturbance information.  Conversely, conditioning on a later event can alter
the inferred distribution of earlier variables, so no-signaling does not
follow from conditioning alone.  Operational no-signaling requires the
remote marginal to be independent of the local intervention after local
outcomes are averaged.  That is a locality condition on the instrument
kernels, consistent with the corrected spatial Bell analysis, not a generic
property of a path measure.

\section{A coherent history is more than an endpoint record}
\label{sec:coherent-history}

The q79--FMO comparison makes the preceding distinction operational. FMO
denotes the Fenna--Matthews--Olson complex. The BEQ calculations considered
here are frozen models of its excitonic and optical response, not biological
measurements and not a derivation of their molecular parameters from MTT.

\subsection{Two projectors with different jobs}\label{sec:route-witness}
For the oriented three-edge route graph, set
\[
B=\begin{pmatrix}-1&-1&0\\0&1&-1\\1&0&1\end{pmatrix},\quad
c=\frac{(-1,1,1)^T}{\sqrt3},\quad P_c=cc^\dagger.
\]
Then \(B^\dagger B/3=I-P_c\). The sign matrix
\(D=\operatorname{diag}(-1,1,1)\) carries the q79 Haar projector to \(P_c\).
A different route-outcome unitary carries the q79 rank-one/rank-two split
to the direct/indirect route split. These are not interchangeable
observables: the cycle projector does not commute with the direct-route
projector. A postselected two-dimensional route-label space cannot contain
an isometric copy of the full three-dimensional carrier \cite{beq02}.

For example, the cycle witness \(\{P_c,I-P_c\}\) gives probabilities
\(1\), \(1/3\), and \(1/9\) on the cycle state, the equal incoherent mixture,
and the all-plus pure state, respectively. Yet the first two states have
the same coarse endpoint populations \((2/3,1/3)\). Population agreement
therefore cannot identify coherent histories. The exploratory lift
\(\sqrt{k}=1/\sqrt{\tau}\) of route rates gives direct weight
\(\tau_{41}/(\tau_{41}+\tau_{42}+\tau_{21})\); its one-third condition is
\(2\tau_{41}=\tau_{42}+\tau_{21}\). This is a declared rate-to-amplitude
diagnostic, not a derived quantum generator or a statistical confidence
interval \cite{beq02,beq03}.

For a positive history operator \(\sigma=\Upsilon\star T_H\), the tester
must preserve the coherent history labels, and its capture
\(p_{\rm cap}=\operatorname{tr}\sigma\) must be positive before conditioning.
Writing \(\rho_H=\sigma/p_{\rm cap}\), witness fidelity
\(\operatorname{tr}(P_c\rho_H)\geq1-\epsilon^2\) implies trace distance at
most \(\epsilon\) from the cycle state, hence at most \(\epsilon\)
total-variation error for every subsequent measurement. Capture and
conditional fidelity are separate experimental quantities; successful
postselection alone establishes neither \cite{beq03}.

\subsection{What the frozen HEOM calculation supplies}\label{sec:heom-history}
The hierarchical equations of motion (HEOM) retain auxiliary density
operators (ADOs) as well as the reduced density matrix. The supplied
seven-site, 300 K reference compares a declared finite set of hierarchy
depths \(D\) and bath-expansion orders \(K\). Its one-time density-matrix
checks are useful baselines but do not determine an intervention-dependent
multi-time process \cite{beq04}.

The later history calculation starts in exciton 4, tags at 50 fs and
100 fs, and retains the ordered histories \((1,1),(2,2),(2,1)\).
Exciton eigenvectors and both tag phases are transported in the same locked
gauge. Applying the tag on both sides of every ADO before subsequent
propagation retains system--bath memory that a reduced-state-only reset
would discard. At the reported \(D6,K1\) rung the capture is approximately
\(0.08027944\), while the conditioned cycle witness is approximately
\(0.29475899\). Thus the ideal cycle prediction and this specified standard
model are separated, not experimentally discriminated \cite{beq05}.

The four-history ADO process slice retains \((1,1),(1,2),(2,1),(2,2)\)
and directly composes 15 CP settings. Its selected three-history block
reproduces the preceding witness and capture, with the reported finite
depth and bath comparisons below their locked tolerance. This resolves the
earlier packet's missing multi-time calculation at this slice, not full
process-tensor tomography. Ground-state bleach (GSB), stimulated emission
(SE), excited-state absorption (ESA), finite pulses, disorder, and
experiment-matched uncertainty are additional types of input
\cite{beq06}. The distinction parallels classical Markovization: an enlarged
state may carry memory exactly while a reduced output alone loses it.

\section{From algebraic readout to a bounded optical instrument}
\label{sec:optical-readout}

\subsection{Dilation and reachability are existence statements}
\label{sec:optical-existence}
For \(K_m=P_1+i^mP_2\) on seven single-exciton levels,
\(\operatorname{rank}(I-K_m^\dagger K_m)=5\), so the minimal unitary
dilation has dimension \(7+5=12\). Ground plus single excitation supplies
only eight levels; the hard-core ground/single/double manifold has
\(1+7+21=29\) and can supply the missing shelving space. The explicit
determinant-corrected dilations give a two-outcome CP instrument. Coherent
deferred selection is valid during the specified number-preserving wait,
not under arbitrary number-changing dynamics \cite{beq07}.

Each of the four additional Hermitian root observables can be written
\(O_a=V_a\Lambda_aV_a^\dagger\). A basis change followed by population
measurement and signed eigenvalue weighting realizes its expectation.
Four settings with seven outcomes describe this direct spectral protocol;
they are not a universal lower bound for all generalized measurements.
The targets are pairwise noncommuting and largely outside the old
exciton-1/2 block. Existing impulsive process-coordinate response arrays
cannot be declared to span these root-state observables without the
preparation/evolution/detection contraction that relates their domains
\cite{beq08}.

The frozen \(K_{ij}=0\), rotating-wave, 29-state control fixture has an
exact strong-regularity and connectivity argument for \(\mathfrak{su}(29)\).
It supplies finite piecewise-control existence with unbounded signed
controls. It does not impose a laboratory amplitude, bandwidth, duration,
or bath tolerance. The separately optimized one-picosecond bounded pulses
pass the strict closed-system gate for only one of four targets. The
weighted-observable searches instead optimize
\(C^\dagger\operatorname{diag}(\lambda)C=O\), but all four returned
candidates miss their stricter gate. Iteration-capped local searches do
not prove unreachable targets \cite{beq09,beq10,beq11}.

\subsection{The bath-active chronology matters}\label{sec:bath-chronology}
The original canonical target-0 five-rung bath ladder fails its declared
comparison tolerance: its largest operator difference is approximately
\(0.00578835>0.005\). A later predeclared \(D3,K1\) higher-corner
comparison stabilizes that waveform and rejects its readout, with relative
residual approximately \(1.00565\). The later result does not retroactively
turn the original five-rung comparison into a pass. The separate weighted
waveform already yields a finite convergence-controlled rejection on its
own five-rung ladder \cite{beq12,beq13,beq15}.

The exact discrete-adjoint result makes bath-aware optimization feasible
without claiming optimal control has been solved. For a finite-rung
segment \(S_j=\exp(\Delta t\,L_j^\dagger)\), its control derivative is a
Fr\'echet derivative of the matrix exponential, evaluated using a block
exponential. Forward and adjoint recursions then differentiate the selected
quadratic readout objective exactly for that discretization. They give no
global minimum or infinite-hierarchy error estimate \cite{beq14}.
Two independently replayed local seeds fail the locked 0.1 relative
operator gate. Convex optimization of terminal-population weights for a
fixed pulse also fails in the nested 7-, 8-, 9-, and 29-outcome families;
only the weight optimization is globally convex. These terminal sectors
are not GSB/SE/ESA pathways. A separate normalized 300 fs, 12-segment search
also fails; its predeclared exit clause stops the proposed 24-segment
continuation for that family \cite{beq16,beq17,beq18}.

\subsection{Full span does not guarantee stable estimation}
\label{sec:bounded-probe}
The molecular-frame finite-probe design first supplies 162 effects of
numerical real Hermitian rank 49. All four targets are in their unbounded
span, but none passes with normalized coefficient norm \(\|w\|_2\leq10\).
With all seven carriers, 882 effects retain full span and three of four
targets pass the unchanged bounded gate. The fourth remains missing; the
three positive results are not erased by that joint failure
\cite{beq19,beq20}.

A predeclared standalone temporal-profile screen then selects the broad
100 fs probe, which passes all four targets at exploratory \(D1,K0\).
The profiles are not pooled. An apparently successful early profile is
rejected by its forward/adjoint integrity check. The broad probe passes
the \(D1,K1\) bath-axis comparison but fails the \(D1,K0\) to \(D2,K0\)
depth comparison; \(D3,K0\) still fails the stacked-effect stability test
\cite{beq21,beq22,beq23}.

The hierarchy-tail diagnostic measures contracting finite increments and
forecasts one further rung. This is not a rigorous tail bound. The subsequent
corrected-accuracy \(D4,K0\) calculation actually supplies that rung:
the maximum per-effect difference is approximately \(0.00039272\), and
the stacked relative difference is approximately \(0.023352\), below
the respective 0.005 and 0.05 thresholds. All four bounded residuals,
approximately \(0.7453,0.5363,0.7359,0.7084\), nevertheless fail. This is a
stable finite-depth negative for this profile, not infinite-HEOM convergence
or a universal optical no-go. Blind depth escalation is not the conclusion
\cite{beq24,beq25}.

\subsection{The physical domain of a nuisance calculation}
\label{sec:physical-tangent}
The complete complex-linear response has 24 rows and 392 ADO coordinates.
After four selected coordinates are separated, 388 nuisance coordinates
span the same rank-18 response image. Six observation combinations remain
nuisance-free, but their selected response rank is zero. Recovery against
one realized nuisance vector therefore does not imply recovery against the
complete complement. Nor does the latter ambient-complex obstruction imply
a physical no-go: Hermiticity, trace, preparation and hierarchy constraints
matter \cite{beq27}.

The later physical-tangent calculation supplies those constraints for the
factorized-root preparation at \(D1,K0\). The trace-free Hermitian root
tangent has 48 real dimensions, rather than the independent 392-real-
coordinate Hermitian ambient hierarchy. Its optical response has reported
rank 18; adjoining the four target rows raises this to 22, exposing four
missing scalar directions. Four stored traceless Hermitian root observables
complete the linear algebra to the stated numerical tolerance. The large
estimator weights and absent calibrated optical realization still matter.
The structural tangent dimension is exact; floating-point row ranks and
residuals are reported source computations, not new exact arithmetic
certificates. Pure-state positivity cones, correlated preparations, and
Hamiltonian/bath parameter tangents require separately chosen domains
\cite{beq26}.

\subsection{Later refinements inside the evolving audit}
\label{sec:evolving-audit}
The candidate audit contains substantive progress after its early snapshot
header. Under unchanged source coordinates and detector norm, adding a
zero-target GSB coordinate cannot lower the existing least-squares residual;
a contractive isotropic transformation cannot enlarge the reachable set.
These restricted statements neither construct the missing full laboratory
response nor rule out a changed physical model \cite{beq01}.

The same audit records distinct external \(K_{ij}\) inputs from the
historical published table and a current-structure reconstruction. Both
alter the signed ESA response, pass exploratory \(D1\) readout, and fail
the matched \(D4\) bounded-readout test. Their effects are genuine model
sensitivity, not MTT parameter selection. On a fixed structure the apparent
21 pair couplings reduce to an external amplitude \(g\) and common twist
\(\phi\); the weight-one dipoles yield only weight-zero and weight-two
pair harmonics. The resulting cone and cross-structure identities permit
conditional source tests. They do not supply an independent covariance
measurement. The separate \(J_{ij}\) ensemble comparison uses ensemble
dispersions, not errors on the means, and is not the double-exciton
\(K_{ij}\) sector \cite{beq01}.

A common ambient circle can carry both sectors with distinct amplitudes
and a fixed relative anchor. A constant 18-degree intertwiner commutes
with circle transport even though 18 degrees is not a \(\mathbb Z_{64}\)
element; interpreting it as such a group element is an additional,
unsupported identification. Charge rows alone do not identify molecular
geometry or that anchor. The frozen B-factor covariance model is a
geometry-uncertainty model, not experimental covariance, and does not resolve
the proposed finite-phase separation \cite{beq01}.

Finally, the audit advances from source-to-generator Lipschitz control to
restricted two-parameter response tangents and analytic remainder bounds.
The \(D1\) tangent has an independent numerical comparison. The executed
\(D4\) tangent and sparse channel Hessian have no interval numerical-error
enclosure or independent \(D4\) validation in the record. Their improved
local radii remain conditional numerical diagnostics, not a certified cover
of the full source family. Detector calibration is also explicit:
\(D_{\rm cal}=RG^T\Sigma^{-1/2}\), with \(R\) the source response,
\(G\) the gain map, and \(\Sigma\) independent positive detector covariance.
Selecting gain or covariance to make the target pass is circular; no
experiment-matched payload for these 882 rows is supplied. These later
refinements are retained without treating packet eligibility flags as
physical evidence \cite{beq01}.

\section{Charged lines, connections, and a conditional repair}
\label{sec:charged-repair}

\subsection{A representation does not select its physical transport}
\label{sec:charged-module}
The frozen 47-atom HF--CIS charge rows
\(q_J=q_{10}\) and \(q_K=q_{11}-q_{00}\) are neutral and linearly independent
over \(\mathbb Q\) as deposited decimals. Tensoring their rank-two span
with a molecular weight-one line gives a charged module with weights
\((1,1)\). A half-turn acts as \(-I_2\); bilinear pair quantities have
weight two and sesquilinear quantities weight zero. Restriction to
\(\mathbb Z_{64}\) is faithful. The exact decimal rank and classifying-stack
pullback do not eliminate electronic-structure or geometry uncertainty
\cite{beq28}.

This construction supplies the charged carrier that the earlier
same-source compatibility packet still listed as missing. A neutral
adjoint/form complex cannot replace it under a nontrivial phase action
\cite{beq47}. The selected nondegenerate seven-site configuration uses
oriented atomic-plane frames and nonzero projected charge vectors. These
frames trivialize the site lines, but the frame-flat connection and the
projected connection \(P\,ds\) are distinct: the latter can have curvature.
A common parallel identification requires trivial relative Hom holonomy,
not merely isomorphic underlying lines \cite{beq29}.

\subsection{Relative phase and two different kinds of loop}
\label{sec:phase-geometry}
On the product of seven unit-frame circles, quotient the diagonal circle
to leave six relative phases. The geometric mechanical connection is
\[
A_F=\frac17\sum_{i=1}^7\alpha_i,\qquad
\beta_i=\alpha_i-A_F,\qquad \sum_i\beta_i=0.
\]
It uses the product frame metric; it is not automatically a nuclear
Eckart connection. The six relative forms need not vanish
\cite{beq30}. On the locked crystal-transfer fixture, the J/K phase
increments lie within an arc shorter than \(\pi\), selecting a shortest
relative-torus path. Its horizontal endpoint comparison is approximately
\(-173.663085\) degrees, or \(6.336915\) after the declared half-turn.
It yields neither the external 18-degree anchor nor an exact
\(\mathbb Z_{64}\) step, and is not itself a q79--FMO physical trajectory
\cite{beq31}.

At fixed configuration the full phase fiber has holonomy
\(\exp(-2\pi i\sum_i n_i/7)\), hence group \(\mu_7\). Its intersection
with \(\mu_{64}\) is trivial, excluding a nontrivial primitive parallel
identification on that full fiber \cite{beq35}. This does not exclude a
geometry-base loop. A synchronous rotation of all seven molecular planes
over a round cap makes the relative forms vanish there. The cap
\(\theta\leq\pi/3\) has area \(\pi\) and boundary holonomy \(-1\), matching
the q79 branch half-turn. The cap normalization is chosen for that
compatibility, not predicted molecular motion. Globally the atom-framed
site lines have \(c_1=0\), so they cannot simultaneously realize a
degree-one full normal sphere with Chern number two \cite{beq42}.

\subsection{Contractible comparison versus a selected endpoint}
\label{sec:interval-descent}
The selected flat q79 family maps to \(B_\nabla U(1)\). Along the
contractible J/K interval, a parallel comparison exists by solving
\(g^{-1}dg=A_q-A_F\), uniquely after one initial fiber map. Its solutions
form a \(U(1)\) torsor. This exact nonempty pullback is not evidence
of a physically selected pairing: interval comparisons exist universally
\cite{beq44}.

The endpoint compiler reduces four future line-comparison rows to two
inputs: a fully selected q79 endpoint and a prospectively selected
q79--FMO pairing rule. It does not supply either input. The endpoint needs
honest descended line/cocycle data, metrics, a common HYM chamber, and
connections; a formal or twisted placeholder is insufficient. The current
hidden projective rank-nine and existential HYM conclusions are retained.
The missing physical visible endpoint and comparison data belong to the
remaining B.HS.01/B.GEO.01 boundary, not to a reopened existential problem
\cite{beq45}.

\subsection{Relative curvature is weaker than equality of connections}
\label{sec:relative-repair}
On the cap the smooth connection is \(A_F=(1-\cos\theta)d\varphi\),
with curvature \(\operatorname{vol}_D\). The punctured q79 representative
is \(A_q=\tfrac12d\varphi\), with coarse current curvature
\(\pi\delta_p\). Their difference
\[
\eta=(\tfrac12-\cos\theta)d\varphi,\qquad
d\eta=\operatorname{vol}_D-\pi\delta_p
\]
has zero tangential boundary value. Thus their relative current class
agrees, while their curvatures and differential characters are not identical.
A ramified/root-stack extension is not an ordinary smooth filled-disk
connection isomorphism \cite{beq38,beq42}.

For the declared round metric and fixed flux \(\pi\), the abelian
Yang--Mills energy is at least \(\pi/2\), with equality exactly at constant
curvature density one. Neumann heat repair
\(F_s=\pi H_N(\kappa s;p,x)\operatorname{vol}_D\) smooths the initial
current for \(s>0\), preserves flux and boundary holonomy, and approaches
that endpoint. Boundary data alone admit infinitely many interior profiles
\(F_F+da\), so they do not select this metric, pairing or law
\cite{beq38}.

On the scalar abelian lane the Maurer--Cartan quadratic bracket vanishes:
the residual is \(dA\). In the relative Coulomb slice its positive squared
residual is the fixed-flux Yang--Mills functional; its gradient gives the
same curvature heat equation. This explains the repair-law shape without
an extra dimensionless coefficient, but \(\kappa>0\) is still an absolute
action/clock scale. This pure positive repair specialization is not the
selected physical signed action and its gauge row is not a physical HYM
moment map \cite{beq41}.

\section{Transported errors and an observable test}
\label{sec:observable-test}

\subsection{Name the metric before bounding the semigroup}
\label{sec:metric-qualification}
For a bounded invertible, domain-, boundary-, and flux-preserving chain
map \(T\), set
\(g_{{\rm tr},k}(x,y)=g_{q,k}(T_k^{-1}x,T_k^{-1}y)\).
Then \(\Delta_{\rm tr}=T\Delta_qT^{-1}\), with exact transported cost and
heat semigroup. Initial isometry is unnecessary. Product preservation
is unnecessary for this scalar curvature action but remains necessary
for the full nonabelian theory \cite{beq40}.

For a separately selected physical metric
\(g_{\rm phys}(x,y)=g_{\rm tr}(x,Ey)\), write
\(C_k=E_{k+1}d_k-d_kE_k\). On a retained bounded sector the exact defect is
\[
\Delta_{\rm phys}-\Delta_{\rm tr}
=d_{k-1}E_{k-1}^{-1}C_{k-1}^{\dagger_{\rm tr}}
 +E_k^{-1}C_k^{\dagger_{\rm tr}}d_k.
\]
The triangle inequality gives a bound \(\delta_k\) from the corresponding
operator norms. Vanishing metric-chain defects suffice for exact physical
intertwining even when \(E\ne I\) \cite{beq40}.

\paragraph{Consumer correction to the frozen error formula.}
The source's coefficient-one estimate \(\kappa s\delta_k\|T\|\) requires
both heat semigroups to be contractions in the \emph{same} norm used for
\(\delta_k\). Self-adjointness in two different metrics does not supply
this hypothesis. In the transported norm, the general Duhamel estimate is
\[
\begin{split}
&\|e^{-\kappa s\Delta_{\rm phys}}T-Te^{-\kappa s\Delta_q}\|\\
&\quad\leq\kappa\delta_k\|T\|
 \int_0^s\|e^{-\kappa(s-u)\Delta_{\rm phys}}\|_{\rm tr}
              \|e^{-\kappa u\Delta_{\rm tr}}\|_{\rm tr}\,du.
\end{split}
\]
For bounded coercive \(E_k\), contraction in the physical metric yields
the valid bound \(K_{E_k}\kappa s\delta_k\|T\|\), where
\(K_{E_k}=\sqrt{\|E_k\|_{\rm tr}\|E_k^{-1}\|_{\rm tr}}\).
Observable evaluation adds \(\|\ell\|\|x\|\) in these same norms.
This qualification preserves the exact transported-metric identities and
the source's commuting diagonal example; it is separately recorded for
source-owner correction, not a blanket downgrade of the result
\cite{beq40}.

Positive-time smoothing cannot manufacture the required bounded inverse.
For an infinite-dimensional compact-resolvent nonnegative Laplacian,
\(e^{-\tau\Delta}\), \(\tau>0\), is compact with dense nonclosed range.
Bounded factors remain compact. On a finite spectral cutoff its inverse
norm grows as \(e^{\tau\lambda_{\max}}\). Therefore a pre-smoothing chain
map or a justified retained finite sector with tail control must come
first; unitary evolution is not excluded by this heat-specific no-go
\cite{beq37}.

\subsection{An exact conditional mode, with a rigorous spectral interval}
\label{sec:cap-observable}
On the round cap, the axisymmetric reducing sector has
\[
L=-\frac{d}{dx}\left((1-x^2)\frac{d}{dx}\right),\quad
\tfrac12\leq x\leq1,
\]
regularity at the pole and Neumann condition \(u'(1/2)=0\).
The least positive root of \(P_\nu'(1/2)=0\) gives
\(\lambda=\nu(\nu+1)\). Normalize its eigenfunction \(\varphi\) in
\(L^2(D)\) and choose its sign positive at the pole. The mean-zero
observable \(\ell_\varphi\) then gives
\(R(s)=\pi\varphi(p)e^{-\kappa\lambda s}\). The flux observable instead
stays constant and cannot determine a clock \cite{beq32}.

The later interval proof excludes earlier roots, uses a pole-free
Dirichlet-to-Neumann monotonicity argument for uniqueness, and bisects an
Arb-certified enclosure. Coarse outward enclosures of its frozen intervals
are
\[
\begin{split}
3.19569115101221&<\nu<3.19569115101222,\\
13.40813308366998&<\lambda<13.40813308367000.
\end{split}
\]
These are imported rigorous intervals, not a new spectral computation
\cite{beq34}. The initial delta current is not an \(L^2\) vector. To apply
the bounded-vector error estimate above, begin at a positive regularization
time or supply a separately justified distribution-space bound. Also
convert the physical \(L^2\) observable norm into the comparison metric
before claiming a numerical error budget.

\subsection{Clock identification and held-out response}
\label{sec:response-identifiability}
For \(Y(t)=C e^{-\beta t}\), unknown nonzero gain and unknown rate are not
identified by one observation. With a positive response, two observations
identify \(\beta=-\log(Y_1/Y_0)/(t_1-t_0)\); a third is held out.
At equal spacing the prediction is \(Y_1^2=Y_0Y_2\). Relative errors
bounded by \(\rho<1\) give a log-defect allowance
\(4[-\log(1-\rho)]\). An additive error needs a positive response lower
bound before using that relative criterion. Fitting the rate is calibration,
not an MTT prediction, and a single exponential is not unique to MTT
\cite{beq33}.

For a positive finite mixture \(Y_n=\sum_j a_jr_j^n\), with
\(a_j>0\) and distinct \(0<r_j<1\),
\[
Y_nY_{n+2}-Y_{n+1}^2
=\sum_{i<j}a_ia_j(r_ir_j)^n(r_i-r_j)^2.
\]
Hence a positive two-by-two Hankel defect excludes a single rate in that
class. The general Hankel rank is the number of distinct modes, not the
three modes in the source's example. Passing a noisy rank-one test does
not prove an exactly single-mode source \cite{beq36}.

Signed and oscillatory responses require the later Prony extension, not
the positivity argument. For distinct nonzero complex nodes and nonzero
amplitudes, the Hankel matrix factors through a Vandermonde matrix and a
diagonal amplitude matrix. A known \(m\)-mode model can be fitted from
\(2m\) training samples and tested on further samples. Signed cancellations
can make a small minor vanish while the full rank is three; conjugate
oscillatory nodes obey a real recurrence. Frequency aliases remain unless
a sampling-band restriction is supplied. If every entry error is at most
\(\delta\), an \(m\)-square Hankel perturbation has norm at most
\(m\delta\); a singular value above that bound certifies retained rank,
but a small singular value alone does not prove exact low rank. Repeated
nodes, continuous spectra, and certified noisy node intervals are further
problems \cite{beq39}.

\subsection{The remaining physical test is not an auxiliary-row count}
\label{sec:physical-gate}
The frozen physical-descent checklist has four supplied rows out of 18:
the charged module, its prospective preselection, the FMO configuration,
and the canonical q79 source base. Its 14 remaining rows ask for an actual
correspondence, parallel connection data, selected amplitudes and anchor,
observable-semigroup/CP-memory and error transport, normalization, clock,
and a prospectively fixed held-out experiment. Audit integrity is not a
physical pass \cite{beq43}.

The later frontier organizes the auxiliary results into finite compatibility,
connection geometry, repair, and response discrimination. A conditional
compiler can close implications within these layers without supplying a
physical input. Its cumulative auxiliary counts cannot be substituted for
the 18 physical obligations. The useful next step is to specify the endpoint
and pairing, then a bounded pre-smoothing chain map and metric defect,
an observable with an error budget, and an independently controlled clock
and detector. None of the preceding algebra requires new biological
measurement principles \cite{beq48}.

\section{The corrected MTT theorem}

\begin{theorem}[Conditional MTT stochastic-shadow theorem]\label{thm:mtt}
Assume a selected MTT branch supplies:
\begin{enumerate}[label=(S\arabic*)]
\item a standard Borel upper state space \(\cZ\) compatible with the
  \(Y^4\times X^6\) branch and its shared-carrier lane data;
\item measurable upper dynamics \(T\), or selected upper transition kernels;
\item a selected probability law \(\rho\) on preparations or captured upper
  states;
\item a measurable observable map \(\pi:\cZ\to Q\).
\end{enumerate}
Then:
\begin{enumerate}[label=(\alph*)]
\item the observed histories possess a unique path law;
\item regular history-dependent conditional kernels exist;
\item failure of the Markov property on \(Q\), or a higher Markov order,
  follows whenever the corresponding conditional-independence witness is
  verified;
\item any action-weight expression gives those kernels only if its action,
  scale, reference measure, and equality to the projected conditional law are
  separately established.
\end{enumerate}
No conclusion about Born probabilities, noncommutative observables, quantum
instruments, process tensors, or CP-divisibility follows from (S1)--(S4)
alone.
\end{theorem}

\begin{proof}
Parts (a) and (b) follow from \Cref{thm:projected-law} and standard
disintegration on standard Borel spaces.  Part (c) follows from
\Cref{def:markov-order,prop:factorization} after a witness has been proved.
Part (d) is \Cref{lem:action-kernel} together with the source distinction in
\Cref{lem:gibbs}.  \Cref{prop:commutative} proves the final limitation.
\end{proof}

\subsection{What is closed and what remains open}

\paragraph{Closed mathematical layer.}
Once a source law or family of kernels is given, the path measure, conditional
kernels, Bayesian posteriors, Markov-order tests, and history-state
Markovization are standard and rigorous.  Projection-induced observed memory
is real and can be demonstrated exactly.

\paragraph{Conditional MTT layer.}
An MTT fixed-point flow and recorder map can instantiate the upper dynamics
and projection.  A selected finite action kernel can instantiate the
conditional law.  These are conditional constructions until the same MTT
branch fixes all their source data.

\paragraph{Open physical layer.}
The decisive missing objects are:
\begin{enumerate}[label=(O\arabic*)]
\item a selected physical capture or preparation measure \(\rho\) for the
  proposed application, beyond the established canonical q79 tier;
\item a derivation of the proposed \(\Delta A\), reference measure, and scale
  from the selected geometry rather than from target probabilities;
\item a proof that no allowed finite MTT state summarizes the claimed memory;
\item a selected family of physical intervention and recorder kernels;
\item a same-source connection between the established noncommutative
  model and the proposed physical preparation, intervention, and detector.
\end{enumerate}

These obligations are testable.  For example, a proposed finite MTT recorder
state can be checked by conditioning two histories with the same recorder
value and comparing their next-step kernels.  A proposed action source can be
tested by computing both the geometric weight and the projected conditional
frequency without fitting one to the other.

\section{Discussion}

\subsection{What survives from the original proposal}

The original physical intuition survives in a narrower and more useful form.
A many-to-one projection can hide upper variables that retain predictive
information.  The lower process can therefore be non-Markovian even when the
upper process is simple.  MTT's fixed-point and bundle language supplies a
natural place for those hidden variables, and the shared carrier offers a
possible common memory channel across projected sectors.

The path-space language is also valuable computationally.  It provides a
precise object to estimate in simulations: conditional kernels indexed by
recorder histories, their total-variation separation, and the minimal
sufficient state needed for prediction.  These are stronger diagnostics than
calling a trajectory ``indivisible'' by inspection.

\subsection{What is no longer claimed}

This paper does not claim that:
\begin{itemize}
\item history dependence eliminates every sequential factorization;
\item failure of one-step Markovianity proves infinite-dimensional memory;
\item a fitted Gibbs action is selected by MTT;
\item a classical event algebra derives the full Hilbert-space quantum
  formalism;
\item measurement is merely conditioning rather than a physical process;
\item no-signaling follows automatically from postselection.
\end{itemize}

Removing those claims strengthens the program.  It leaves a theorem that can
be used without ambiguity and identifies exactly which future calculation
would turn a flexible stochastic encoding into an MTT prediction.

\subsection{Relation to established work}

Chains with complete connections provide a mature framework for processes
whose conditional probabilities depend continuously on a long or infinite
past \cite{fernandez2005,johansson2003}.  Their uniqueness and mixing theorems
require specific regularity hypotheses; ``summable variations'' is not a
synonym for arbitrary history dependence.  Those results can be applied after
an MTT kernel has been selected and its variation bounds have been verified.

Quantum non-Markovianity is richer because interventions can disturb the
system and the relevant objects are noncommutative multi-time maps.  Process
tensors and CP-divisibility are therefore appropriate comparison targets, not
automatic consequences of the classical shadow construction.  The revised
MTT path is:
\[
\begin{array}{c}
\text{selected upper source and dynamics}\\
\downarrow\\
\text{projected classical path law and memory diagnostics}\\
\downarrow\quad\text{(additional noncommutative source theorem)}\\
\text{quantum process tensor or channel family}.
\end{array}
\]

\section{Conclusion}

MTT can rigorously support a history-dependent stochastic shadow, but only
after a probability source has been selected.  The resulting lower process
may be non-Markovian and even have infinite Markov order on the recorded state
space.  This is a substantive consequence of projection.  It is nevertheless
state-relative: the process is Markovian on its history space and may admit a
small sufficient-state compression.

The corrected result therefore does two things at once.  It preserves the
useful MTT insight that hidden upper coherence can appear as lower memory, and
it prevents that insight from being mistaken for a derivation of quantum
probability or absolute indivisibility. The new q79--FMO account makes that
boundary concrete. Canonical operational quantum structure, coherent-history
calculations, and exact charged-line compatibility are already available.
Finite numerical readout profiles and conditional geometric repair do not
yet select a physical molecular correspondence. The next research target
is that same-source preparation and intervention model, with calibrated
observable transport, followed by a test of which memory remains on the
chosen recorder state.

\small
\bibliographystyle{plain}

% BEGIN MTT MANAGED COMPUTATIONAL EVIDENCE
\section*{Computational Evidence and Reproducibility}
The classical path-law proofs are independent of the numerical packets.
The contextual sections cite all 48 assigned BEQ frozen artifacts at commit
\texttt{f141a20e a23c5c3f f19cc216 1c0e226e 29ade8a7} of the
\href{https://github.com/PeterNero/mtt-results-repro}{curated results repository}.
The frozen manifest is \path{release/result_manifest.json}; its SHA-256 is
\begin{quote}\small\ttfamily
5715b0da 6a54df29 d43a53fe 4b9ae932\\
bf6d4cd9 fc43e8c0 ff2e2254 56df13e3
\end{quote}
Each bibliography link identifies its exact artifact bytes. Local metadata
and the review fragment bind result identifiers to source hashes and this
manuscript. The evolving candidate audit is read through its later sections;
its preceding-snapshot self-reference and older eligibility headers are
historical context, not current authority over refined results.

This revision performs source reading, byte-hash validation, and small
bounded algebraic consistency checks. It does not replay HEOM, control
optimization, interval spectral isolation, or molecular calculations, and
it does not infer independent verification from packet booleans. Negative
readout results retain their exact fixture, norm, and finite-rung scope.
Conditional geometry retains its metric, domain, pairing and clock inputs.
Current canonical operational q79 and finite mathematical conclusions are
not reopened by historical source-selection language. No imported result
changes theorem ownership or promotes the physical q79--FMO gate.
% END MTT MANAGED COMPUTATIONAL EVIDENCE

\bibliography{main}

\end{document}
